% GENERATED FILE. Edit canonical lessons, then run npm run generate:books.
% canonical-source-hash: fnv1a64:22cd391bd6d72045
% canonical-lessons: TA-C539-pullankulal, TA-C539-vinai, TA-C539-vaci, TA-C539-mittu, TA-C539-icaiyamai

\chapter{Music Lessons}
\label{ch:ta-music-lessons}
\hlchaptermodality{539}

\begin{chapteropening}
\textbf{By the end of this chapter:} \emph{I can talk about playing an instrument.}

Complete the last lesson of chapter 539: I can talk about playing an instrument.
\end{chapteropening}

\section[\texorpdfstring{\ta{புல்லாங்குழல்} (pullāṅkuḻal)}{புல்லாங்குழல் (pullāṅkuḻal)}]{\ta{புல்லாங்குழல்} (pullāṅkuḻal) — a flute}

\label{lesson:TA-C539-pullankulal}



\begin{quote}
Before the new one: say the Tamil for a loom, then the Tamil for silk.
\end{quote}

\subsection*{You'll want to know: \ta{புல்லாங்குழல்}}
\textbf{\ta{புல்லாங்குழல்}} — \emph{pullāṅkuḻal} — \textquotedblleft{}a flute\textquotedblright{}.

\begin{quote}
\ta{இடையன்} \ta{மரத்தடியில்} \ta{புல்லாங்குழல்} \ta{ஊதினான்}.
\end{quote}

\begin{quote}
\emph{iṭaiyaṉ marattaṭiyil pullāṅkuḻal ūtiṉāṉ.}
\end{quote}

\begin{quote}
The herdsman played the flute under a tree.
\end{quote}

A flute is blown, so it takes \ta{ஊது} (\emph{ūtu}): \ta{புல்லாங்குழல்} \ta{ஊது}.

\begin{cousinweb}
\textbf{\ta{புல்லாங்குழல்}} (\emph{pullāṅkuḻal}) ends in the native \ta{குழல்} (\emph{kuḻal}), \textquotedblleft{}a tube, a flute, a pipe\textquotedblright{} (DEDR 1818); the origin of the first part is not settled.
\end{cousinweb}

\begin{grammarlens}[title={What you've built}]
One more word for describing what you lived and felt.
\end{grammarlens}

\subsection*{Your turn}
\begin{itemize}
\raggedright
  \item \emph{Say it:} \emph{pullāṅkuḻal}
  \item \emph{Say it:} \emph{pullāṅkuḻal}, once more
  \item \emph{Say it:} \emph{iṭaiyaṉ marattaṭiyil pullāṅkuḻal ūtiṉāṉ.} — The herdsman played the flute under a tree.
  \item \emph{Recall:} say the Tamil for a loom, then the Tamil for silk, then say \emph{pullāṅkuḻal} again
\end{itemize}

\begin{tcolorbox}[breakable,colback=teal!4,colframe=teal!35!black,title={Before you move on}]
What does \ta{புல்லாங்குழல்} mean? (\textquotedblleft{}A flute\textquotedblright{}.) Say it once more.
\end{tcolorbox}
\section[\texorpdfstring{\ta{வீணை} (vīṇai)}{வீணை (vīṇai)}]{\ta{வீணை} (vīṇai) — a veena}

\label{lesson:TA-C539-vinai}



\begin{quote}
Before the new one: say the Tamil for silk, then the Tamil for a flute.
\end{quote}

\subsection*{You'll want to know: \ta{வீணை}}
\textbf{\ta{வீணை}} — \emph{vīṇai} — \textquotedblleft{}a veena\textquotedblright{}.

\begin{quote}
\ta{அவள்} \ta{பத்து} \ta{வயது} \ta{முதல்} \ta{வீணை} \ta{கற்றுக்கொள்கிறாள்}.
\end{quote}

\begin{quote}
\emph{avaḷ pattu vayatu mutal vīṇai kaṟṟukkoḷkiṟāḷ.}
\end{quote}

\begin{quote}
She has been learning the veena since she was ten.
\end{quote}

A large stringed instrument of South Indian classical music.

\begin{cousinweb}
\textbf{\ta{வீணை}} (\emph{vīṇai}) is borrowed from Sanskrit \emph{vīṇā}, \textquotedblleft{}a lute\textquotedblright{} (CDIAL 12048); Tamil writes the long final vowel as -\ta{ஐ}.
\end{cousinweb}

\begin{grammarlens}[title={What you've built}]
One more word for describing what you lived and felt.
\end{grammarlens}

\subsection*{Your turn}
\begin{itemize}
\raggedright
  \item \emph{Say it:} \emph{vīṇai}
  \item \emph{Say it:} \emph{vīṇai}, once more
  \item \emph{Say it:} \emph{avaḷ pattu vayatu mutal vīṇai kaṟṟukkoḷkiṟāḷ.} — She has been learning the veena since she was ten.
  \item \emph{Recall:} say the Tamil for silk, then the Tamil for a flute, then say \emph{vīṇai} again
\end{itemize}

\begin{tcolorbox}[breakable,colback=teal!4,colframe=teal!35!black,title={Before you move on}]
What does \ta{வீணை} mean? (\textquotedblleft{}A veena\textquotedblright{}.) Say it once more.
\end{tcolorbox}
\section[\texorpdfstring{\ta{வாசி} (vāci)}{வாசி (vāci)}]{\ta{வாசி} (vāci) — to play an instrument; to read aloud}

\label{lesson:TA-C539-vaci}



\begin{quote}
Before the new one: say the Tamil for a flute, then the Tamil for a veena.
\end{quote}

\subsection*{You'll want to know: \ta{வாசி}}
\textbf{\ta{வாசி}} — \emph{vāci} — \textquotedblleft{}to play an instrument; to read aloud\textquotedblright{}.

\begin{quote}
\ta{அவர்} \ta{கச்சேரியில்} \ta{வீணை} \ta{வாசித்தார்}.
\end{quote}

\begin{quote}
\emph{avar kaccēriyil vīṇai vācittār.}
\end{quote}

\begin{quote}
He played the veena at the concert.
\end{quote}

Past (I …): \ta{வாசித்தேன்}, \emph{vācittēṉ}.

Any instrument can take \ta{வாசி}; a flute may also take \ta{ஊது} (\emph{ūtu}), to blow. It also means \textquotedblleft{}to read\textquotedblright{}.

\begin{cousinweb}
The origin of \textbf{\ta{வாசி}} (\emph{vāci}) is not settled; one suggestion links it with Sanskrit \emph{vādya}, \textquotedblleft{}a musical instrument, to be played\textquotedblright{} (CDIAL 11511).
\end{cousinweb}

\begin{grammarlens}[title={What you've built}]
One more word for describing what you lived and felt.
\end{grammarlens}

\subsection*{Your turn}
\begin{itemize}
\raggedright
  \item \emph{Say it:} \emph{vāci}
  \item \emph{Say it:} \emph{vāci}, once more
  \item \emph{Say it:} \emph{avar kaccēriyil vīṇai vācittār.} — He played the veena at the concert.
  \item \emph{Recall:} say the Tamil for a flute, then the Tamil for a veena, then say \emph{vāci} again
\end{itemize}

\begin{tcolorbox}[breakable,colback=teal!4,colframe=teal!35!black,title={Before you move on}]
What does \ta{வாசி} mean? (\textquotedblleft{}To play an instrument; to read aloud\textquotedblright{}.) Say it once more.
\end{tcolorbox}
\section[\texorpdfstring{\ta{மீட்டு} (mīṭṭu)}{மீட்டு (mīṭṭu)}]{\ta{மீட்டு} (mīṭṭu) — to pluck strings}

\label{lesson:TA-C539-mittu}



\begin{quote}
Before the new one: say the Tamil for a veena, then the Tamil for to play an instrument.
\end{quote}

\subsection*{You'll want to know: \ta{மீட்டு}}
\textbf{\ta{மீட்டு}} — \emph{mīṭṭu} — \textquotedblleft{}to pluck strings\textquotedblright{}.

\begin{quote}
\ta{அவள்} \ta{வீணையின்} \ta{கம்பிகளை} \ta{மெதுவாக} \ta{மீட்டினாள்}.
\end{quote}

\begin{quote}
\emph{avaḷ vīṇaiyiṉ kampikaḷai metuvāka mīṭṭiṉāḷ.}
\end{quote}

\begin{quote}
She gently plucked the strings of the veena.
\end{quote}

Past (I …): \ta{மீட்டினேன்}, \emph{mīṭṭiṉēṉ}.

The strings take -\ta{ஐ}: \ta{கம்பிகளை} \ta{மீட்டு}.

\begin{cousinweb}
\textbf{\ta{மீட்டு}} (\emph{mīṭṭu}) is native Dravidian: DEDR 4848 gives \ta{மீட்டு} (\emph{mīṭṭu}), \textquotedblleft{}to fillip the strings of a lute, to play the lute\textquotedblright{}.
\end{cousinweb}

\begin{grammarlens}[title={What you've built}]
One more word for describing what you lived and felt.
\end{grammarlens}

\subsection*{Your turn}
\begin{itemize}
\raggedright
  \item \emph{Say it:} \emph{mīṭṭu}
  \item \emph{Say it:} \emph{mīṭṭu}, once more
  \item \emph{Say it:} \emph{avaḷ vīṇaiyiṉ kampikaḷai metuvāka mīṭṭiṉāḷ.} — She gently plucked the strings of the veena.
  \item \emph{Recall:} say the Tamil for a veena, then the Tamil for to play an instrument, then say \emph{mīṭṭu} again
\end{itemize}

\begin{tcolorbox}[breakable,colback=teal!4,colframe=teal!35!black,title={Before you move on}]
What does \ta{மீட்டு} mean? (\textquotedblleft{}To pluck strings\textquotedblright{}.) Say it once more.
\end{tcolorbox}
\section[\texorpdfstring{\ta{இசையமை} (icaiyamai)}{இசையமை (icaiyamai)}]{\ta{இசையமை} (icaiyamai) — to compose music}

\label{lesson:TA-C539-icaiyamai}



\begin{quote}
Before the new one: say the Tamil for to play an instrument, then the Tamil for to pluck strings.
\end{quote}

\subsection*{You'll want to know: \ta{இசையமை}}
\textbf{\ta{இசையமை}} — \emph{icaiyamai} — \textquotedblleft{}to compose music\textquotedblright{}.

\begin{quote}
\ta{அவர்} \ta{இந்தத்} \ta{திரைப்படத்துக்கு} \ta{இசையமைத்தார்}.
\end{quote}

\begin{quote}
\emph{avar intat tiraippaṭattukku icaiyamaittār.}
\end{quote}

\begin{quote}
He composed the music for this film.
\end{quote}

Past (I …): \ta{இசையமைத்தேன்}, \emph{icaiyamaittēṉ}.

The film or song takes -\ta{க்கு}: \ta{படத்துக்கு} \ta{இசையமை}. A composer is an \ta{இசையமைப்பாளர்}.

\begin{cousinweb}
\textbf{\ta{இசையமை}} (\emph{icaiyamai}) is a native compound: \ta{இசை} (\emph{icai}), \textquotedblleft{}music\textquotedblright{}, and \ta{அமை} (\emph{amai}), \textquotedblleft{}to fit together, to arrange\textquotedblright{} (DEDR 162).
\end{cousinweb}

\begin{grammarlens}[title={What you've built}]
One more word for describing what you lived and felt.
\end{grammarlens}

\subsection*{Your turn}
\begin{itemize}
\raggedright
  \item \emph{Say it:} \emph{icaiyamai}
  \item \emph{Say it:} \emph{icaiyamai}, once more
  \item \emph{Say it:} \emph{avar intat tiraippaṭattukku icaiyamaittār.} — He composed the music for this film.
  \item \emph{Recall:} say the Tamil for to play an instrument, then the Tamil for to pluck strings, then say \emph{icaiyamai} again
\end{itemize}

\begin{tcolorbox}[breakable,colback=teal!4,colframe=teal!35!black,title={Before you move on}]
What does \ta{இசையமை} mean? (\textquotedblleft{}To compose music\textquotedblright{}.) Say it once more.
\end{tcolorbox}
